\documentclass[10pt,a4paper]{amsart}
\usepackage[margin=19mm]{geometry}
\usepackage{amsmath,amssymb,mathtools}
\usepackage[scr=rsfs,cal=euler]{mathalfa}
\usepackage{xfrac}
\usepackage[unicode,hidelinks,bookmarks=false]{hyperref}
\newcommand{\ie}{i.e.\ }
\newcommand{\cf}{cf.\ }
\newcommand{\resp}{respectively\ }
\newcommand{\Subsection}{\subsection}
\providecommand{\nobreakdash}{\nobreak\hbox{-}}
\setlength{\emergencystretch}{2em}
\multlinegap0pt
\usepackage{algorithm}
\usepackage{algpseudocode}
\usepackage{enumerate}
\usepackage{subcaption}
\usepackage{comment}
\usepackage{dsfont}
\algnewcommand\algorithmicinput{\textbf{Input:}}
\algnewcommand\Input{\item[\algorithmicinput]}
\algnewcommand\algorithmicoutput{\textbf{Output:}}
\algnewcommand\Output{\item[\algorithmicoutput]}
\newtheorem{assumption}{Assumption}
\newtheorem{proposition}{Proposition}
\newtheorem{remark}{Remark}
\newtheorem{theorem}{Theorem}
\newcommand{\tp}{^\top}
\newcommand{\inv}{^{-1}}
\newcommand{\R}{\mathbb{R}}
\renewcommand{\S}{\mathcal{S}}
\newcommand{\simulator}{\S}
\newcommand{\ham}{\mathcal{H}}
\newcommand{\half}{\frac{1}{2}}
\DeclareMathOperator*{\argmin}{arg\,min}
\newcommand{\ip}[2]{\langle #1, #2 \rangle}
\newcommand{\normal}{\mathcal{N}}
\newcommand{\uns}{\mathrm{uns}}
\newcommand{\ud}{\mathrm{d}}
\newcommand{\G}{\mathcal{G}}
\newcommand{\grad}{\nabla}
\newcommand{\todo}[1]{{\color{blue} #1}}
\newcommand{\backward}[1]{\overset{\leftarrow}{#1}}
\newcommand{\one}{\mathds{1}}
\newcommand{\diag}{\mathrm{diag}}
\newcommand{\id}{\mathcal{I}}
\renewcommand{\L}{\mathcal{L}}
\newcommand{\cb}{\color{blue}}
\title[A dual ensemble Kalman filter approach to robust control]{A Dual Ensemble Kalman Filter Approach to Robust Control of Nonlinear Systems: An Application to Partial Differential Equations: the robust-control proposition for the nonlinear model}
\author[Anant A. Joshi, Saviz Mowlavi, Mouhacine Benosman]{Anant A. Joshi, Saviz Mowlavi, Mouhacine Benosman}
\date{Source: arXiv:2508.21684v1 (CC BY 4.0)}
\begin{document}
\begin{abstract}
This paper considers the problem of data-driven robust control design for nonlinear systems, for instance, obtained when discretizing nonlinear partial differential equations (PDEs). A robust learning control approach is developed for nonlinear affine in control systems based on Lyapunov redesign technique. The robust control is developed as a sum of an optimal learning control which stabilizes the system in absence of disturbances, and an additive Lyapunov-based robustification term which handles the effects of disturbances. The dual ensemble Kalman filter (dual EnKF) algorithm is utilized in the optimal control design methodology. A simulation study is done on the heat equation and Burgers partial differential equation.
\end{abstract}
\maketitle

\noindent\emph{Source and licence.} A Dual Ensemble Kalman Filter Approach to Robust Control of Nonlinear Systems: An Application to Partial Differential Equations. Version arXiv:2508.21684v1 (CC BY 4.0), distributed under the Creative Commons Attribution 4.0 International licence, http://creativecommons.org/licenses/by/4.0/, as recorded in the arXiv OAI licence metadata for this exact version and shown on the abstract page https://arxiv.org/abs/2508.21684v1. Full rights notice: licenses/arxiv-2508.21684-CC-BY-4.0.txt.\par\smallskip

\noindent\textit{Editorial scope.} Counted unit: the paper's robust-control proposition for the nonlinear affine-in-control system, printed as the second of the paper's two Propositions (source0.tex lines 440--443), delivered together with the paper's complete original proof (source0.tex lines 444--457) as the single primary proof body. No mathematics is written, repaired or re-derived: statement and proof are the paper's own text line for line, in the paper's own notation ($\R$, $\tp$, $\inv$, $\grad$, $\bar{u}$, $u_d$, $\lambda$, $b^{\dagger}$, $V$, $\half$). Required context retained uncounted: the whole of Section 1 (source0.tex lines 116--168) with the definition of the controlled system, the notation paragraph and the roadmap; the whole of Section 2 (lines 169--370) with the linear reduced-order simulator, the optimal control problem, the paper's first Proposition and its complete proof, and the data-driven simulator-based implementation; the whole of Section 3 (lines 371--718) with the nonlinear discretized model, its stabilizing control, the assumptions on $b$ and on the disturbance bound, the counted proposition, its proof, the data-driven implementation, Section 4 on the application to forced nonlinear partial differential equations, and the future-work section; and Appendix A with the dual ensemble Kalman filter algorithms and the simulation details (lines 725--805), so that every section, proposition, assumption, remark, numbered display, algorithm, figure caption and citation of the delivered range is retained and every cross-reference of the delivered text resolves inside the delivered document. The delivered text carries the paper's own citation commands over exactly the thirty-two bibliography keys of the paper, and the thirty-two entries are transcribed verbatim from the paper's own BibTeX product kept by the prep lane in eprint/main.bbl, in that file's own order, so the printed bibliography numbers of the delivered unit are the paper's own (1 to 32). Reference adaptation: none was needed and none was made, and no reference command of the delivered text was rewritten. Printed slips kept literal: the paper's own forms are delivered as printed, among them the misspelling methoology in the Introduction, the broken sentence reading and they are $\chi_j$ is the indicator functions of in the simulation details, the paper's own use of the word pseudo inverse, its unlabelled propositions, and its habit of writing the robust term as $u_d$ throughout. Layout and class: the paper's own class is ieeeconf in conference mode on letter paper with the IEEE two-column style file, which is NOT redistributed with this package and is not used to build it; the delivered wrapper is the lane's pinned amsart editorial wrapper at 10pt on A4, so the two-column IEEE layout, its fonts, its margin overrides and its running heads are replaced by a single-column editorial layout. The wrapper reproduces the paper's own macro block (its transpose, inverse, R, S, Hamiltonian, half, argmin, inner-product, normal, uns, differential, G, gradient, backward, indicator, diagonal, identity and L macros, together with the algorithm input and output redefinitions), the paper's own four theorem-like declarations for assumption, proposition, remark and theorem, and the packages the retained text requires (algorithm, algpseudocode, enumerate, subcaption and comment). The paper's own preamble loads the algpseudocodex package for its algorithm blocks, whose dependency fifo-stack is not installed in this runtime; the two retained algorithm environments are therefore set with the algorithmicx package algpseudocode, which provides every command the paper's algorithms use (Input, State, If, ElsIf, EndIf, For, EndFor and Return) and prints the same statements in the same order, with a typographic difference in the end-markers and line labels of the algorithm blocks. Apart from that package substitution, the wrapper reproduces the paper's own macro block (its transpose, inverse, R, S, Hamiltonian, half, argmin, inner-product, normal, uns, differential, G, gradient, backward, indicator, diagonal, identity and L macros, together with the algorithm input and output redefinitions) and the paper's own four theorem-like declarations for assumption, proposition, remark and theorem; that is a page-layout and algorithm-typesetting difference of the wrapper only, and no formula, symbol, hypothesis or argument changes. Assets: the paper contains eight graphics-inclusion commands in five figure environments and no PDF-page inclusion, input or include directive. Six of those eight commands lie inside the retained span 458--718, inside the paper's three remaining figure environments whose captions and labels are retained verbatim; exactly those six graphics-inclusion commands for the files under the archive folder Fig were omitted from the delivered text, so the delivered document contains no graphics-inclusion command at all, no image file is copied into or redistributed with this package, and the package declares no asset dependency. Those six omissions are recorded in delivery-evidence.json as a bounded source comparison with a fidelity receipt that binds the delivered bytes, the source bytes and the item hash; nothing else inside that span was altered, and the two remaining graphics-inclusion commands of the paper lie in the omitted toy-example sections. Omitted from this package and not redistributed: the paper's preamble (lines 1--105), its front matter (lines 106--115: the IEEE-style title, the author block with the footnote carrying the affiliations and the abstract, of which the title, the author names and the abstract are reproduced in the wrapper title block), the bibliography commands (lines 720--721), and the two closing sections of simulation plots and toy examples (lines 808--904). No publisher class, no publisher style file, no font and no upstream script is redistributed. A full rights notice is delivered at licenses/arxiv-2508.21684-CC-BY-4.0.txt.
\section{Introduction}
In this paper, we are primarily interested in the robust control of nonlinear affine in control systems modelled as 
\begin{align}
\dot{x}(t) = a(x(t)) + b(x(t))u(t) + d(t,x), \quad x(0) = x_0
\label{eq:control-system}
\end{align}
where $x(t) \in \R^n$, $u(t) \in \R^m$ and $d(t,x) \in \R^l$ for all $(t,x)$, and $d$ is regarded as an unknown disturbance bounded in norm by a known real valued function $\lambda(t,x)$. The emphasis is on systems where we do not have explicit access to $a$ and $b$, but we can access trajectories of the system, either generated by a simulator or collected from real-life experiments. For simplicity, we focus in the remaining of the paper on the case of simulated trajectories. 
Furthermore, we target the application of our algorithms to the control of PDEs, when their discretized model (\ref{eq:control-system}) is available as a simulator. 

Indeed, PDE control is challenging because a PDE is by nature infinite dimensional and may have strong nonlinearities. 
The reduce-then-design approach is a well researched control methoology, which involves discretization of the PDE followed by dimensionality reduction, which yields a model amenable to the application of standard model based control approaches \cite{ %hinze2005proper, 
leibfritz2007numerical, hovland2008explicit, barbagallo2009closed,sipp2016linear, %gao2017active, 
tsolovikos2020estimation}. There is a recent body of work using data driven methods for building more accurate reduced order models of the PDE \cite{%lee2020model,
bhattacharya2021model, fresca2021comprehensive, kaiser2021data}. While more accurate reduced order models are beneficial to the performance of the controller, their complexity makes them computationally challenging to implement. 

% In this work therefore, we focus more on the design of the controller, given either a coarse reduced order model, or we work with the full space-discretized nonlinear model. The emphasis is on building a controller only given a simulator for the model, without access to the actual model parameters. 
In this work, we focus more on the design of the controller with emphasis on obtaining the control policy with only simulator for the model, without access to the actual model parameters. See for example \cite{duriez2017machine,blanchard2021bayesian,fan2020reinforcement, Garnieretal21} for some other recent efforts in obtaining data-driven controllers using a simulator of the PDE.   
In particular, we work on the stabilization problem, which aims to drive the PDE state to zero, by posing it as an optimal control problem. The first step is to discretize the PDE in space, to yield a high-dimensional nonlinear system of ordinary differential equations (ODEs). We consider two different cases for the simulator availability: 
\begin{enumerate}
    \item a linear reduced order simulator, obtained using Dynamic Mode Decomposition with control (DMDc) \cite{proctor2016dynamic,zhang-2024} is available.
    \item an extension to the case where a simulator for the full high-dimensional nonlinear discretized model of the PDE is available.
\end{enumerate}

We build on previous work in this area \cite{zhang-2024}  by taking a robust learning control approach. We construct the robust control as a sum of two parts: an optimal learning controller, which we expect to produce stabilization in the absence of disturbances, and an additional term that builds on the optimal control term and is based on the Lyapunov redesign approach \cite{khalil} to suppress the effect caused by disturbances. 
%The central component of our design recipe is the designing the optimal control, for which we use the dual ensemble Kalman filter (dual EnKF) algorithm \cite{joshi-2022}. The dual EnKF algorithm simulates multiple interacting copies of the system model, to yield an approximation to the optimal control. 
The optimal control is approximated using the dual ensemble Kalman filter (dual EnKF) algorithm \cite{joshi-2022}, which simulates multiple interacting copies of the system.

The ensemble Kalman filter (EnKF) \cite{yang-2016,amir-2019} has historically been an algorithm used for filtering, and is a key numerical method especially for high-dimensional systems, for example in weather prediction \cite{evensen2006} (see \cite{bishop-2023} for more references). It features the design of an interacting particle system to sample from the posterior density of a filtering problem to provide a state estimate. 
The dual EnKF algorithm, inspired from the EnKF, computes the optimal control by converting the control problem into the problem of sampling from an appropriate probability density, through the log-transform duality between optimal control and filtering \cite{fleming-1982}. Such an approach for solving optimal control by posing it as a sampling problem is well studied
(see \cite{todorov-2007,kappen-2005,rawlik2013stochastic,hoffman-2017,levine-2018}). 
% The dual EnKF  leverages the fundamental duality between optimal control and filtering \cite{fleming-1982}. 
% The optimal control problem is converted to the problem of sampling from a density. 
% The dual EnKF converts the optimal control problem into the problem of sampling from a probability density, through the use of log transform duality between optimal control and filtering \cite{fleming-1982}.
The novelty of the dual EnKF lies in the design of an interacting particle system, inspired from the EnKF, to solve the sampling problem. 
% When applied to sampling, the EnKF exhibits a distinct advantage in its performance on systems of high dimension \cite{evensen2006}. 
The dual EnKF-based control exhibits a distinct advantage in its performance on systems of high dimension -- a trait inherited from the EnKF. It performs almost two orders of magnitude faster when compared with policy gradient type approaches \cite{mihailo-2021-tac,fazel-2018-mlr} as was computationally demonstrated in \cite{joshi-2022}. 


The paper is organized as follows. In Section \ref{sec:linear} we introduce and solve the problem for the case when the system is linear time invariant (LTI) to clarify the main ideas. Then we present the extension to the nonlinear affine in control case in Section \ref{sec:nonlinear}. 
%In Section \ref{sec:numerical} we present some simple numerical examples to test and validate the algorithm on. 
Finally, in Section \ref{sec:pde} we present an application to the stabilization of PDEs, with the heat equation and Burgers' equation as examples. 

% \noindent \textbf{Related literature:}

% \noindent \textbf{Contributions:}

% Emphasis on the simulator based setting, enkf does really well in high dimensions, comparison in \cite{joshi-2022}

% give some idea of dual EnKF

\textbf{Notation:} %$\one$ represents a vector of ones, $\normal(0,1)$ represents the normal distribution with mean 0 and variance 1, 
$|\cdot|$ denotes the Euclidean norm, $\tp$ denotes the transpose of a matrix, and $\id$ refers to the identity matrix.

\section{Solution for linear reduced order simulator}
\label{sec:linear}
%THE STABILIZING CONTROLLER SHOULD BE FIRST AND THEN ROBUSTIFICATION IN SECTION b

Consider the simplified case when \eqref{eq:control-system} is a linear time invariant (LTI) system:
\begin{align}
\dot{x}(t) = Ax(t) + Bu(t) + d(t,x)
\label{eq:lin-control-system}
\end{align}
where as earlier, $x(t) \in \R^n$, $u(t) \in \R^m$ and 
$d(t,x) \in \R^l$ for all $(t,x)$, and $d$ is regarded as a disturbance. From here on we may suppress the $t$ argument to improve readability. We construct the robust control as a sum of a stabilizing control and an additional robustification term. We first present a recipe to obtain the stailizing control, and then to obtain the robust control, and lastly a simulator based method to obtain the two. 

\subsection{Stabilizing control}


Consider, for any arbitrary but fixed $T > 0$, the optimal control problem
\begin{subequations}
\begin{align}
% \min_{u(\cdot)} \mathcal{J}(u(\cdot)) & \coloneqq \nonumber \\  x(T) \tp G x(T) \, + \,  & \half\int_{0}^{T} |Cx(t)|^2 + u(t)\tp R u(t) \ud t \\
\min_{u(\cdot)} & \bigg( x(T) \tp G x(T) + \half\int_{0}^{T} \underbrace{|Cx(t)|^2 + u(t)\tp R u(t)}_{=: \L(x(t),u(t))} \ud t \bigg)  \\
\text{ s.t. } & \text{system }\eqref{eq:lin-control-system} \text{ with zero disturbance, that is, } d \equiv 0
\end{align}
\label{eq:LQR}
\end{subequations}
where $C \in \R^{n_1 \times n}$ (for arbitrary $n_1 > 0$), $G \in \R^{n \times n}$ and $R \in \R^{m \times m}$, and  define $Q := C\tp C$. 
\begin{assumption}
We make the following assumptions about the structure of the optimal control problem \eqref{eq:LQR}
\begin{enumerate}[(i)]
\item $(A,B)$ is controllable and $(A,C)$ is observable
\item $R,G \succ 0$
\end{enumerate}

\end{assumption}
Then there exists a positive definite solution $\{P(t) : t \in [0,T]\}$ to the differential Riccati equation (DRE) \cite[Chapter 3]{k-sivan}:
\begin{align*}
-\dot{P} = A\tp P + PA - PBR\inv B\tp P + Q, \quad P_T = G
\end{align*}
which converges to $\bar{P} \succ 0$ as $T \to \infty$ and $\bar{P}$ solves the algebraic Riccati equation (ARE)
\begin{align*}
0 = A\tp \bar{P} + \bar{P}A - \bar{P}BR\inv B\tp \bar{P} + Q.
\end{align*}
Moreover, the control $u = -\bar{K}x$ with $\bar{K} := R\inv B\tp \bar{P}$ makes the system asymptotically stable \cite[Theorem 3.7]{k-sivan}. 

\subsection{Robust control}

Suppose that for the system \eqref{eq:lin-control-system} in the case of no disturbance, there exists a stabilizing control $u = -\bar{K}x$ and a strictly positive definite Lyapunov function $V(x) = \half x\tp \bar{P} x$ with $\dot{V} = -x\tp \bar{P}(A - B\bar{K})x \le 0$ along the controlled trajectories with zero disturbance. We let $\bar{K}$ be the gain obtained from the optimal control demonstrated previously and $\bar{P}$ be the solution of the ARE.  
%A constructive recipe to obtain both will be specified in the subsequent section.

We are interested in the idea of disturbance rejection using Lyapunov redesign \cite[Chapter 14]{khalil}. The idea is to design a robust control $u_d$ such that the controller $u = -\bar{K}x + u_d$ makes the system \eqref{eq:control-system} asymptotically stable in the presence of disturbance. To that end, we make the following assumption, and then present a design methodology for $u_d$. %\todo{edit}

\begin{assumption}
\begin{enumerate}[(i)]
    \item The rank of $B$ is $n$. 
    \item There exists a known $\lambda$ such that $0 \le |d(t,x)| < \lambda(t,x) < \infty$ for each $(t,x) \in [0,\infty) \times \R^n$.
\end{enumerate}
\label{assn:B}
\end{assumption}
\begin{remark} 
    Assumption \ref{assn:B}-(i) is required to motivate the theoretical derivation for the linear system, but we relax it in our implementations for PDE control. 
\end{remark}
% \begin{remark}
% If the rank of $B$ is at least $n$, for every $y \in \R^n$ there exists at least one solution $u_y$ to the linear equation $Bu = y$ given by $u_y = (B\tp B)\inv B\tp y$. 
% \end{remark}
% \begin{remark}
% Assumption 2 implies that the dimension of controls is at least as large as the dimension of states.
% \end{remark}



Consider the controller $u = -\bar{K}x + u_d$ with 
% \begin{align*}
%     u_d := - \frac{\lambda(t,x)}{| \bar{P}x|} v, \quad v := B^{\dagger} \bar{P}x, \quad B^{\dagger} := (B\tp B)\inv B\tp,
% \end{align*}
\begin{align*}
    u_d := - \frac{\lambda(t,x)}{| \bar{P}x|} B^{\dagger} \bar{P}x, \quad B^{\dagger} := (B\tp B)\inv B\tp,
\end{align*}
where $B^{\dagger}$ denotes the Moore-Penrose pseudoinverse. 
%Under Assumption \ref{assn:B}, we have that $Bv = \bar{P}x$ hence $Bu_d = - \frac{\lambda}{| \bar{P}x|} \bar{P}x$.
The following result demonstrates the effectiveness of the robust control.
% $u_d := - \frac{\lambda(t,x)}{| \bar{P}x|} v$ with $v := B^{\dagger} \bar{P}x$ and $B^{\dagger} := (B\tp B)\inv B\tp$. 
% Then $Bv$ is the orthogonal projecion of $\bar{P}x$ onto the column span of $B$. 
% 
% So
% \begin{align*}
% \ip{\bar{P}x}{Bv} = \ip{Bv}{Bv} = \left| B B^{\dagger} \bar{P}x\right|^2 \le | BB^{\dagger} |^2 \cdot |\bar{P}x |^2 \\
% \implies \ip{Bu_d}{\bar{P}x} = -\frac{\lambda}{|\bar{P}x|} \left| B B^{\dagger} \bar{P}x\right|^2.
% \end{align*}
% 


% We summarise these calculations in the following result. 
\begin{proposition}
Using the control $u = -\bar{K}x + u_d$ renders the system \eqref{eq:control-system} asymptotically stable. 
% \label{prop:robust}
\end{proposition}
\begin{proof}
We follow the method in \cite[Chapter 14.2]{khalil}, using the same Lyapunov function $V(x) = \frac{1}{2} x\tp \bar{P}x$ as before. 
The quadratic form $V$ is strictly positive definite (by assumption) hence is a valid Lyapunov function. Taking derivative along system trajectories, 
\begin{align*}
\dot{V}(x) &= -x\tp \bar{P}(A - B\bar{K})x  +  (Bu_d + d)\tp \bar{P} x \\
& \le  | d | \cdot | \bar{P}x| -\lambda |\bar{P} x| \\
& \le  (|d| - \lambda ) \, | \bar{P}x| < 0.
\end{align*}
% Moreover, $u = -\bar{K}x$ by assumption is an asymptotically stable controller for \eqref{eq:lin-control-system} when there is no disturbance \cite[Theorem 3.7]{k-sivan} and therefore $\dot{V} \le 0$. 
%With $D \ne 0$, the choice of $u_d$ renders $\dot{V} < 0$, hence makes the system asymptotically stable. 
For the first inequality we recall that $-x\tp \bar{P}(A - B\bar{K})x \le 0$ and use Cauchy-Schwarz inequality. 
Next by properties of the Moore-Penrose pseudoinverse, $BB^{\dagger}\bar{P}x$ is the orthogonal projection of $\bar{P}x$ onto the column span of $B$. 
Hence, under Assumption \ref{assn:B}-(i), we have $BB^{\dagger}\bar{P}x = \bar{P}x$ therefore $Bu_d = - \frac{\lambda}{| \bar{P}x|} \bar{P}x$.
\end{proof}
 \begin{remark}
     When implementing $u_d$ we add a regularizing parameter to avoid division by zero, which makes the system asymptotically stable till it enters a ball around the origin, i.e., practical stability. 
 \end{remark}

\subsection{Data-driven (Simulator-based) implementation}
In this section, we present a method to implement the control $u = -\bar{K}x + u_d$ with only access to a disturbance-free system simulator of \eqref{eq:control-system}. We use the dual ensemble Kalman filter (dual EnKF) algorithm \cite{joshi-2022} for the same. To use the algorithm, we need the following assumptions:
\begin{assumption}
We have access to the following:
\begin{enumerate}[(i)]
\item Knowledge of optimal control matrices $Q,R,G$. %\todo{Q needed?}
\item Simulator of the dynamical system with no disturbance, that is, we have access to function evaluations of $\S(x,u) := Ax + Bu$. Moreover, we assume we can run the simulator backward in time, that is, to find a trajectory of the system by specifying the terminal condition. 
\end{enumerate}
\label{assn:access}
\end{assumption}
\begin{remark}
\label{rmk:evaluate}
With access to a perfect simulator, one may exactly find the model matrices $A$ and $B$ in $n+m$ evaluations of the simulator (set $u=0$ and evaluate the simulator at the basis vectors of $\R^n$ to find $A$ and equivalently for $B$). However, the utility of the simulator is revealed when we consider the nonlinear case, especially for high-dimensional systems like partial differential equations, where estimating the state dynamics in this manner is not possible. 
\end{remark}


% \subsubsection{Obtain controller using simulator:}
The emphasis is on obtaining the controller in a model-free way. The following three steps are done (which are elaborated upon after listing them):
\begin{enumerate}
\item Find an approximation to $\bar{P}$, the solution of the ARE, using the dual EnKF algorithm \cite{joshi-2022}. See Appendix \ref{app:linenkf} for details
\item Find an approximation $\bar{u}^{(N)}$ for $\bar{u} := -\bar{K}x$ using \cite[Algorithm 2]{joshi-2022} (recalled in Appendix \ref{app:ubar})
\item Find an approximation $u_d^{(N)}$ for $u_d$ using Algorithm \ref{alg:ud}
\end{enumerate}

\textbf{Step 1:}
Using the simulator, we find $\bar{P}^{(N)}$, an approximation to the solution of the ARE $\bar{P}$ by running the dual EnKF algorithm of \cite{joshi-2022}, which simulates $N$ copies of the dynamical system  along with a mean-field coupling term to approximate the solution of the ARE. The algorithm is provided in Appendix \ref{app:linenkf}. 
% This algorithm needs the following:
% \begin{enumerate}
% \item Access to simulator $\S$ (in particular, ability to run the simulator backward in time). 
% \item Access to function evaluations of $x \mapsto Cx$
% \item Explicit knowledge of $R$
% \end{enumerate}

\textbf{Step 2:} %\todo{add some inspiration for Ham}
To help evaluate the optimal control in a model-free way, define the Hamiltonian,
\begin{align}
\ham(x,u) &:= (\bar{P}^{(N)}x)\tp(Ax + Bu) + \half(x\tp Q x + u\tp R u) \nonumber
\\ &= (\bar{P}^{(N)}x)\tp \S(x,u) + \half\L(x,u) \label{eq:ham-lin}
\end{align}
The Hamiltonian is constructed so that it can be evaluated using function calls of $\S$. For a fixed $x$, it is a quadratic function of $u$ with the unique minima at the optimal control. 
%Since the Hamiltonian relies on $x,\bar{P}^{(N)}$ and function calls of $\S$ and $c$, it can be evaluated using the simulator. 
Therefore, $\bar{u} = \argmin_{u} \ham(x,u)$ and the minimization can be carried out using gradient estimation as shown in \cite[Algorithm 2]{joshi-2022} (recalled in Appendix \ref{app:ubar}) or using zero-order methods, such as \cite{bach-2016}. 

\textbf{Step 3:}
Similarly, to find $u_d$ we solve $u_d = -\lambda \argmin_{u} | Bu - \frac{\bar{P}x}{| \bar{P}x|} |$ in Algorithm \ref{alg:ud}.  If $B$ is known, one may directly use the Moore-Penrose pseudo inverse. If $B$ is unknown, one may use zero order optimization methods \cite{bach-2016} where $Bu = \S(0,u)$ can be obtained using only access to simulator or one may estimate $B$ as specified in Remark \ref{rmk:evaluate} and use the pseudo inverse. 
{Using zero order optimization is  preferred over estimating $B$ in cases when the simulator is very high dimensional. }
%Algorithm \ref{alg:ud} displays a gradient estimation type method. One may use zero order optimization methods \cite{bach-2016} where $Bu = \S(0,u)$ can be obtained using only access to the simulator.

\begin{algorithm}
\begin{algorithmic}[1]
\Input System state $x$, regularizing parameter $r$, robust gain $\lambda$, $\bar{P}^{(N)}$
\State $r_1 := \max( \, |\bar{P}^{(N)}x| \, , \, r \, ) $
\If{$B$ is known}
\State $B^{\dagger} := (B\tp B)\inv B\tp$
\State $v^{(N)} := r_1\inv B^{\dagger} \bar{P}^{(N)}x$
\ElsIf{$B$ is unknown}
\State $v^{(N)} := \argmin_{v} | \simulator(0,v) - \frac{\bar{P}^{(N)}x}{r_1} |$
\EndIf
\State \Return $u^{(N)}_d := -\lambda v^{(N)}$
\end{algorithmic}
\caption{Algorithm to find $u_d$}
\label{alg:ud}
\end{algorithm}


% \begin{algorithm}
% {\cb 
% \begin{algorithmic}[1]
% \Input System state $x$, regularizing parameter $r$, robust gain $\lambda$, $\bar{P}^{(N)}$
% \State $r_1 := \max( \, |\bar{P}^{(N)}x| \, , \, r \, ) $
% \If{$B$ is known}
% \State $B^{\dagger} := (B\tp B)\inv B\tp$
% \State $v^{(N)} := r_1\inv B^{\dagger} \bar{P}^{(N)}x$
% \ElsIf{$B$ is unknown}
%     \If{Assumption \ref{assn:B}-(i) holds and it is easy to estimate $B$}
%         \State estimate $B$ using Remark \ref{rmk:evaluate}
%         \State Go to line 2
%     \Else
%         \State $v^{(N)} := \argmin_{v} | \simulator(0,v) - \frac{\bar{P}^{(N)}x}{r_1} |$ using \cite{bach-2016}
%     \EndIf
% \EndIf
% \State \Return $u^{(N)}_d := -\lambda v^{(N)}$
% \end{algorithmic}
% }
% \caption{Algorithm to find $u_d$}
% \label{alg:ud}
% \end{algorithm}


\section{Extension for nonlinear discretized model of the PDE}
\label{sec:nonlinear}

Let us consider now the main result of this paper, dealing with the case of the nonlinear affine in control system \eqref{eq:control-system}. We will extend the robust control design methodology presented in Section \ref{sec:linear} to \eqref{eq:control-system}. Similar to the previous section, we will obtain the robust control as a sum of a stabilizing control and a robustification term, and then present a simulator based methodology to obtain both terms. 
% \begin{align}
% \dot{x}(t) = a(x(t)) + b(x(t))u(t) + d(t,x)
% \label{eq:control-system}
% \end{align}
% where $x(t) \in \R^n$, $u(t) \in \R^m$ and 
% $d(t,x) \in \R^l$ for all $(t,x)$, and $d$ is regarded as a disturbance. 
% % Hereon we may suppress the $t$ argument to improve readability. 
% We are interested in the idea of disturbance rejection using Lyapunov redesign \cite[Chapter 14]{khalil}. The idea is to design a robust control $u_d$ such that the controller $u = -\bar{K}x + u_d$ makes the system \eqref{eq:control-system} asymptotically stable for non-zero disturbance, that is, for non-zero $D$ and non-zero $d$. To that end, we make the following assumption, and then present a design methodology for $u_d$. \todo{edit}
\subsection{Stabilizing control}
In an effort to find a stabilizing control $\bar{u}$, we consider, for any arbitrary but fixed $T > 0$, the optimal control problem
\begin{subequations}\label{eq:nLQR}
\begin{align}
% \min_{u(\cdot)} \mathcal{J}(u(\cdot)) \coloneqq & \mathcal{G}(x(T)) + \half \int_{0}^{T} c(x(t)) + u(t)\tp R u(t) \ud t  \\
\min_{u(\cdot)} & \bigg( \mathcal{G}(x(T)) + \half \int_{0}^{T} \underbrace{c(x(t)) + u(t)\tp R u(t)}_{=: \L(x(t),u(t))} \ud t \bigg) \\
\text{ s.t. } & \text{system }\eqref{eq:control-system} \text{ with zero disturbance, that is, } d \equiv 0
\end{align}
\end{subequations}
where $c,\G$ are non-negative real valued function, and $R \in \R^{m \times m}$ is symmetric and strictly positive definite. The value function for the problem is defined as the cost-to-go,
\begin{align*}
    \phi(s,x) &:= \min_{u(\cdot)}  \bigg( \mathcal{G}(x(T)) + \half \int_{s}^{T} \L(x(t),u(t)) \ud t \bigg) \\
\text{ s.t. } & \text{system }\eqref{eq:control-system} \text{ with } X_s = x \text{ and } d \equiv 0
\end{align*}
and it satisfies the Hamilton Jacobi Bellman partial differential equation. The optimal control is computed as $\bar{u}(t,x) = -R\inv b(x)\tp \grad \phi(t,x)$ \cite{liberzon}. 
% \begin{remark} (MB. I do not agree here as well, since we do have a stabilizing Lyapunov function, we can use the results of Lyapunov inverse optimality to find the corresponding cost function that we are optimizing via this controller. I would suggest to remove this Remark.)
% While it is in general not possible to prove convergence properties of the optimal control law thus obtained, the optimal control formulation is used in applications which require stabilization of nonlinear systems, see for example, \cite{tedrake-notes,zhang-2024,hovland2008explicit,rawlik2013stochastic}.
% \end{remark}
\subsection{Robust control}
\begin{assumption}
\begin{enumerate}[(i)]
    \item There exists a control law $\bar{u}$ which makes \eqref{eq:control-system} with zero disturbance asymptotically stable. Moreover, there exists a strictly positive definite Lyapunov function $V$ with $\dot{V} < 0$ along the controlled trajectories with zero disturbance. 
    \item The rank of $b(x)$ is $n$ for all $x \in \R^n$. 
    \item There exists a known $\lambda$ such that $0 \le |d(t,x)| < \lambda(t,x) < \infty$ for each $(t,x) \in [0,\infty) \times \R^n$.
\end{enumerate}
\label{assn:b}
\end{assumption}
\begin{remark}
    Assumption \ref{assn:b}-(ii) is required for the theoretical result, but will be relaxed in the PDE control implementation. 
\end{remark}

% \begin{remark}
% If the rank of $B$ is at least $n$, for every $y \in \R^n$ there exists at least one solution $u_y$ to the linear equation $Bu = y$ given by $u_y = (B\tp B)\inv B\tp y$. 
% \end{remark}
% \begin{remark}
% Assumption \ref{assn:B} implies that the dimension of controls is at least as large as the dimenision of states.
% \end{remark}

Similar to the linear case, consider the controller $u = \bar{u} + u_d$ with 
\begin{align*}
    u_d := - \frac{\lambda}{| \grad V |} b^{\dagger} \grad V, \quad b^{\dagger} := (b\tp b)\inv b\tp.
\end{align*}
% $u_d := - \frac{\lambda(t,x)}{| \grad V(x) |} v$ with $v := b^{\dagger} \grad V(x)$ and $b^{\dagger} := (b\tp b)\inv b\tp$. 
% Then $bv$ is the orthogonal projection of $\grad V(x)$ onto the column span of $b$. So
% \begin{align*}
% \ip{\grad V(x)}{bv} = \ip{bv}{bv} = \left| b b^{\dagger} \grad V(x) \right|^2 \le | bb^{\dagger} |^2 \cdot |\grad V(x) |^2 \\
% \implies \ip{bu_d}{\grad V(x)} = -\frac{\lambda}{|\grad V(x)|} \left| b b^{\dagger} \grad V(x)\right|^2.
% \end{align*}

% Under Assumption \ref{assn:B}, we have $bv = \grad V(x)$ hence $Bu_d = - \frac{\lambda}{| \bar{P}x|} \bar{P}x$. For the same Lyapunov function $V(x) = x\tp \bar{P}x$ we have,
% \begin{align*}
% \dot{V}(x) &= \grad V (x(t)) \tp (a(x(t)) + b(x(t))\bar{u}(t) + b(x(t))u_d(t) + Dd(t,x)) \\
% & \le  (| d | \cdot | \grad V(x)| -\lambda |\grad V (x)|) \\
% & \le (|d| - \lambda ) |\grad V (x)| < 0.
% \end{align*}

% We summarise these calculations in the following result. 

\begin{proposition}
Using the control $u = \bar{u} + u_d$ renders the system \eqref{eq:control-system} asymptotically stable. 
% \label{prop:robust}
\end{proposition}

\begin{proof}
We follow the method in \cite[Chapter 14.2]{khalil}. 
Taking the derivative of $V$ along system trajectories, 
\begin{align*}
\dot{V}(x) &= \grad V (x(t)) \tp \bigg( a(x(t)) + b(x(t))\bar{u}(t) \\ & \quad + b(x(t))u_d(t) + d(t,x) \bigg) \\
& \le  (| d | \cdot | \grad V(x)| -\lambda |\grad V (x)|) \\
& \le (|d| - \lambda ) |\grad V (x)| < 0.
\end{align*}
% Moreover, $u = -\bar{K}x$ by assumption is an asymptotically stable controller for \eqref{eq:lin-control-system} when there is no disturbance \cite[Theorem 3.7]{k-sivan} and therefore $\dot{V} \le 0$. 
%With $D \ne 0$, the choice of $u_d$ renders $\dot{V} < 0$, hence makes the system asymptotically stable. 
For the first equality we recall Assumption \ref{assn:b} and Cauchy-Schwarz inequality. 
%Next, by properties of the Moore-Penrose pseudoinverse, $bv$ is the orthogonal projection of $\grad V(x)$ onto the column span of $b$. Thus under Assumption \ref{assn:b}, we have that $bv = \grad V(x)$ therefore $bu_d = - \frac{\lambda}{| \grad V(x)|} \grad V(x)$.
Moreover, under Assumption \ref{assn:b}-(ii), $bu_d = - \frac{\lambda}{| \grad V(x)|} \grad V(x)$.
\end{proof}



\subsection{Data-driven (Simulator-based) implementation}

To approximate $\bar{u}$ and $u_d$ we use an approach similar in spirit to Section \ref{sec:linear}. First, we use the nonlinear dual EnKF algorithm \cite{joshi-2022} to approximate the gradient of value function $\grad \phi^{(N)}(x)$. Implementation details can be found in Appendix \ref{app:nlenkf}.
Then we define the nonlinear counterpart of the Hamiltonian  
\begin{align} \label{eq:ham-nl}
\ham(x,u) &:= (\grad \phi(x)^{(N)})\tp \S(x,u) + \half\L(x,u)
\end{align}
where the simulator is now nonlinear, that is, $\S(x,u) = a(x) + b(x)u$. The Hamiltonian can again be evaluated using function calls of the simulator. Moreover, it is quadratic in the control, and can be minimized using \cite[Algorithm 2]{joshi-2022} (recalled in Appendix \ref{app:ubar}), or zero order optimization approaches \cite{bach-2016}. Similar to the linear control case, to find $u_d$ we use $u_d = -\lambda \argmin_{u} | b(x)u - \frac{\grad V(x)}{| \grad V(x)|} |$ as given in Algorithm \ref{alg:ud} (replacing $\bar{P}^{(N)}x$ by $\grad V^{(N)}(x)$) where if $b$ is not known, the optimization can be done by zero-order approaches \cite{bach-2016} or by estimating $b$ similar to Remark \ref{rmk:evaluate}.

\section{Application to forced nonlinear PDEs}
\label{sec:pde}
We consider PDEs of the form
\begin{align*}
    \frac{\partial z}{\partial t}(t,y) + \mathcal{F}(z(t,y)) = \omega(t,y),
\end{align*}
where $\mathcal{F}$ is a differential operator that specifies the structure of the PDE, $z$ is the state of the PDE and $\omega$ is the external input. The functions $z, \omega : \R_+ \times [0,L] \to \R$. Mathematically, the goal of stabilization is to make 
\begin{align}\label{eq:L2objective}
\lim_{t \to \infty}\|z(t,\cdot)\|_{L^2} := \lim_{t \to \infty}\int_{0}^{L} |z(t,y)|^2 \ud y = 0.    
\end{align}
For numerical implementation, we consider a time interval of $[0,T]$ and discretize the PDE in space on a grid of $p$ uniformly-spaced points in $[0,L]$ to obtain $z_p(t) \in \R^p$ for each $t \ge 0$, so that the PDE is reduced to a set of $p$ ODEs. We choose $m$ basis functions $\{\chi_j\}_{j=1}^{m}$ to discretize the control as 
$\omega(x,t) = \sum_{j=1}^m \chi_j(x) U_j(t)$ where $\chi_j : [0,L] \to \R$ is the indicator function of $[\frac{j-1}{m},\frac{j}{m}]$ and $U_j : [0,T] \to \R$. 
Then $U := (U_1,U_2,\ldots,U_m)$ is interpreted as the control. The discretized PDE has nonlinear affine in control form,
\begin{align}
\frac{\ud z_p}{\ud t} + F(z_p) = BU
\label{eq:discrete-pde}
\end{align}
where $F$ approximates the derivatives and $B$ obtained by discretizing $\{\chi_j\}_{j=1}^{m}$. To study the effect of disturbances, we let $U(t) = u(t) + d(t)$ where $u(t)$ is the control action applied to the system and $Bd(t)$ is the effective disturbance acting on the system. We consider here disturbances added directly to the control, which is a special case of the theory presented earlier. In the simulations, we consider only time varying disturbances, so we denote disturbance as $d(t)$. 
We use the Controlgym library in Python \cite{zhang2023controlgym} for numerical implementation. 
Now we first recall the heat equation and Burgers equations, give some information about the implementation of both PDEs, then go on to discuss the simulation results.
\subsection{Heat equation}

The heat equation is given by 
\begin{align*}
    \frac{\partial z}{\partial t}(t,y) - \nu \frac{\partial^2 z}{\partial y^2}(t,y) = \omega(t,y).
\end{align*}
The value of $\nu = 0.002$ is used. Since the PDE is linear, upon discretization, \eqref{eq:discrete-pde} reduces to an LTI system, hence we use the robust control method design discussed in Section \ref{sec:linear} to implement the robust control. The dimension of the control used is $m=10$, while the dimension of the state is $n=100$, thus relaxing Assumption \ref{assn:B}-(i) in the implementation.  Other simulations details can be obtained in Appendix \ref{app:heat}.
\begin{figure}
\centering
    \begin{subfigure}[b]{0.45\textwidth}
        \centering
        
        \caption{}
        \label{fig:heatmap-heat}
    \end{subfigure}
%
    \begin{subfigure}[b]{0.45\textwidth}
        \centering
        
        \caption{}
        \label{fig:time-heat}
    \end{subfigure}
\caption{Results for control of heat equation. (a) Mean of $\frac{\|z(T,\cdot)\|_{L^2}}{\|z(0,\cdot)\|_{L^2}}$ over 100 different simulations (b) Mean and variance of $\|z(t,\cdot)\|_{L^2}$ for $d_0=0.1, \lambda=0.2$. }
\label{fig:heat}
\end{figure}

% \begin{figure}
% \centering
%     \begin{subfigure}[b]{0.2\textwidth}
%         \centering
%         \includegraphics[scale = 0.25]{Fig/heatmap-heat.pdf}
%         \caption{tt}
%         \label{fig:heatmap-heat}
%     \end{subfigure}
% %
%     \begin{subfigure}[b]{0.2\textwidth}
%         \centering
%         \includegraphics[scale = 0.25]{Fig/time-heat.pdf}
%         \caption{tt}
%         \label{fig:time-heat}
%     \end{subfigure}
% \caption{hwfk}
% \label{fig:heat}
% \end{figure}

\subsection{Burgers' equation}

% The goal is to stabilize the Burgers equation in the presence of disturbances using the methodology developed so far. 
Recall the Burgers' equation, 
\begin{align*}
\frac{\partial z}{\partial t}(t,y) + z(t,y) \frac{\partial z}{\partial y}(t,y) - \nu \frac{\partial^2 z}{\partial y^2}(t,y) = \omega(t,y).
\end{align*}
% with $z, \omega : [0,T] \times [0,L] \to \R$. \todo{add definition and boundary conditions}. Mathematically, the goal is to make 
% \begin{align*}
% \lim_{t \to \infty}\|z(t,\cdot)\|_{L^2} := \int |z(t,x)|^2 \ud x = 0.    
% \end{align*}

% For numerical implementation, we discretize the PDE in space on a grid of $p$ uniformly spaced points in $[0,L]$ to obtain $z_p(t) \in \R^p$ for each $t \ge 0$, so that the PDE is reduced to a set of $p$ ODEs. We choose $m$ basis functions $\{\chi_j\}_{j=1}^{m}$ to discretize the control as 
% $\omega(x,t) = \sum_{j=1}^m \chi_j(x) U_j(t)$ where $\chi_j : [0,L] \to \R$ and $U_j : [0,T] \to \R$. 
% Then $U := (U_1,U_2,\ldots,U_m)$ is interpreted as the control. The discretized PDE has the nonlinear affine in control form,
% \begin{align}
% \frac{\ud z_p}{\ud t} + F(z_p) = B(z_p)u.
% \label{eq:discrete-pde}
% \end{align}
% where $F$ approximates the derivatives and $B$ obtained by discretizing $\{\chi_j\}_{j=1}^{m}$. To study the effect of disturbance, we let $U(t) = u(t) + d(t)$ where  where $u$ is the control action applied to the system and $d$ is a disturbance acting on the system. We refer the reader to \cite{controlgym} for the implementation details. 

% \todo{write about control discretization and disturbance modelling. To study robust control, we assume that the action is perturbed by a disturbance, that is, $U(x,t) = u(x,t) + d(x,t)$ where $u$ is the action which is intended to be applied and $d$ is a disturbance acting on the system.}

We present two sets of results to demonstrate the performance of the robust control algorithm developed in the paper when applied to the Burgers equation. In the first set of results, we calculate the robust control using a simulator of the linear reduced order model of the PDE, and in the second set of results, the robust control is calculated using a simulator of the full nonlinear discretized PDE. In both results, the obtained control is tested by applying it on the full nonlinear discretized PDE. We study both sets of results for two values of viscosity, $\nu = 0.02,0.002$.


\noindent \textbf{Control design using linear reduced model:} 
We perform a model reduction on the Burgers' equation to obtain a linear system using the Dynamic Mode Decomposition for control (DMDc) algorithm \cite{zhang-2024}. The DMDc yields a projection matrix $\Phi \in \R^{n \times p}$ with orthogonal rows, where $n$ is the dimension of the reduced state (usually $n \ll p$), and an LTI system  
\begin{align}
\dot{x} = Ax + Bu
\label{eq:DMDc}
\end{align}
where $x \in \R^n$ is the reduced state and satisfies the relation $x = \Phi z_p$, and $u \in \R^m$ where as described, $n$ and $m$ are chosen by the user. The discretized PDE state can be reconstructed as $z_p = \Phi\tp x$. DMDc yields a discrete-time system, but we transform it into its continuous-time equivalent, since we design control in the continuous time domain. 

Using a simulator for the obtained linear system \eqref{eq:DMDc}, we get an approximation to the optimal stabilizing LQ controller, with the linear EnKF methodology described in Section \ref{sec:linear}. 
%The LQ control is computed using the reduced LTI model at each point in time, but applied to the full nonlinear PDE. 
That is, we choose optimal control weights $Q,R,G$ (as described in \eqref{eq:LQR}) and find $\bar{P}^{(N)}$ for the LTI system \eqref{eq:DMDc}. The stabilizing control is calculated as $u = -R\inv B\tp \bar{P} x$ where $x = \Phi z_p$. We use the robust control methodology described in Section \ref{sec:linear} to design $u_d$. 

% \subsubsection{Control design using reduced model}

% We perform a model reduction on the Burgers equation to obtain a linear system using the Dynamic Mode Decomposition for control (DMDc) algorithm \cite{zhang-2024}. The DMDc yields a projection matrix $\Phi \in \R^{n \times p}$ where $n$ is the dimension of the reduced state (usually $n << p$) and an LTI system  
% \begin{align}
% \dot{x} = Ax + Bu
% \label{eq:DMDc}
% \end{align}
% where $x \in \R^n$ is the reduced state and satisfies the relation $x = \Phi z_p$, and $u \in \R^m$ where as described, $n$ and $m$ are chosen by the user. 

% Using a simulator for the obtained linear system \eqref{eq:DMDc}, we get an approximation to the optimal stabilizing LQ controller, with the linear EnKF methodology described in Section \ref{sec:linear}. 
% %The LQ control is computed using the reduced LTI model at each point in time, but applied to the full nonlinear PDE. 
% That is, we choose optimal control weights $Q,R,G$ (as described in \eqref{eq:LQR}) and find $\bar{P}^{(N)}$ for the LTI system \eqref{eq:DMDc}. The stabilizing control is calculated as $u = -R\inv B\tp \bar{P} x$ where $x = \Phi z_p$. 

% We study the effect of two types of disturbances, sinusoidal with $d(t) := d_0sin(t)$ and constant $d(t) := d_0$. We use the same robust control methodology as described in Section \ref{sec:linear} to design $u_d$. The obtained control is applied to the space discretized full nonlinear PDE \eqref{eq:discrete-pde} and the observed behaviour of the PDE is illustrated in Figure \ref{fig:}. The experiments are performed for two values of viscosity, $\nu = 0.02,0.002$. \todo{what is exactly plotted? also include discussion on the results, see comments below}

% \todo{Mention number of particles, dimension of state and control, value of $Q,R,G$}

% We study two types of disturbances, of $d$ that is $d(t) := d_0sin(t)$ (in Figure \ref{fig:dmdc-sint}) and $d(t) := d_0$ (in Figure \ref{fig:dmdc-const}). We use the same robust control methodology as described in Proposition \ref{prop:robust}, and study the effect of the robust control on the full PDE for various values of $\nu$ (as the value of $\nu$ increases, the PDE becomes less non-linear). In each figure, the average of the norm of the initial condition is also given as a reference to show that stabilisation has indeed ocurred. 

% For $\nu = 0.1, 0.01 $, there is a very clear trend in the results. For $d_0 = 0$, $\lambda = 0.02$ works the best and for other values of $d_0$, we see that $\lambda = 20d_0$ works best. This holds for both sinusoidal and constant disturbances. 

% For $\nu = 0.001$ however, the case with $d_0 = 0,0.01$ are problematic since for all values of $\lambda > 0$ which we tried, the system gets de-stabilized due to the robust control. It may not be very clearly visible in the colour plots (since the colours may look similar making it hard to draw this inference) but when one looks at the actual numerical values, one sees that there is indeed de-stabilization. For $d_0 = 0.1,1$ however, the $\lambda = 20d_0$ trend is restored.

% \subsubsection{Control design using full nonlinear model}

\noindent \textbf{Control design using full nonlinear model:}
Using a simulator for the full nonlinear PDE \eqref{eq:discrete-pde}, we obtain a stabilizing control $\bar{u}$ and robust controller $u_d$ with the nonlinear dual EnKF methodology discussed in Section \ref{sec:nonlinear}. We choose $c(x) := |x|^2$ and $\mathcal{G}(x) = |x|^2$ in \eqref{eq:nLQR}. 
%In the simulator for the full nonlinear PDE \eqref{eq:discrete-pde}, 
The dimension of the state used is $p=128$ and the dimension of the control is $m=10$, thus relaxing Assumption \ref{assn:b}-(ii) in the implementation. 
% We study the effect of the same two types of disturbances, sinusoidal with $d(t) := d_0sin(t)$ and constant $d(t) := d_0$.

% The obtained control is tested on full nonlinear PDE \eqref{eq:discrete-pde} 
% and the results are in Figure \ref{fig:heatmap-full}. The experiments are performed 
% for two values of viscosity, $\nu = 0.02,0.002$. 

% \noindent \textbf{Discussion of results:}
\subsection{Discussion of results}
We study the effect of two types of disturbances, sinusoidal with $d(t) := d_0\sin(t)$ and constant $d(t) := d_0$. The $\lambda$ function is chosen as a constant function. The obtained control is applied to the full discretized nonlinear PDE \eqref{eq:discrete-pde}. % and the observed behaviour of the PDE is illustrated in Figure \ref{fig:}. 
% The experiments are performed for two values of viscosity, $\nu = 0.02,0.002$. 
Both PDEs are initialized for 100 randomly sampled iid initial conditions (inspired from the previous work \cite{zhang-2024}): 
%$z( t=0,x) = \alpha \, \mathrm{sech}(\frac{1}{\beta}(x - \frac{1}{2L}))$ with $L=1$, $\alpha \sim \mathrm{unif}(0.9, 1.1)$ and $\beta \sim \mathrm{unif}(0.04, 0.06)$. 
\begin{align*}
    z(0,x) &= \alpha \, \mathrm{sech}\left(\frac{1}{\beta}\left({x} - \frac{1}{2L}\right)\right), \quad L=1, \; \\ 
    \alpha &\sim \mathrm{unif}(0.9, 1.1), \quad \beta \sim \mathrm{unif}(0.04, 0.06).
\end{align*}
The 100 trajectories are simulated for the case of zero control (that means $\bar{u} = u_d = 0$), and then with the robust control given by various values of $\lambda$ (the stabilizing control $\bar{u}$ is obtained from the dual EnKF). These simulations are repeated for various values of $\lambda$, $d_0$, both types of disturbances mentioned, and the values of $\nu$ mentioned. 

% For both PDEs, we make two plots. The first plot is a heat map depicting the mean value (over the 100 simulations) of $\\|z(T,\cdot)\|_{L^2}$ -- the energy of the state at the final time. Recall that the control objective is to drive $\|z(T,\cdot)\|_{L^2}$  to zero. Moreover, in a second plot, we show the mean and variance (of the 100 simulations) of the $\|z(t,\cdot)\|_{L^2}$ versus $t$ for $d_0 = 0.1$ and $\lambda = 0.2$. The trajectory is generated using three different control policies -- ``uncontrolled" trajectories are when the control is zero, ``optimal controlled" trajectories have $\lambda = 0$ (that is, only the optimal stabilizing control is applied and robust control is zero) and ``robust controlled" trajectories are with $\lambda = 0.2$. The plots for heat equation, Burgers with DMDc and Burgers with full nonlinear model are found in Figures \ref{fig:heat}, \ref{fig:dmdc}, and \ref{fig:full}, respectively. Other simulations details can be obtained in Appendix \ref{app:burgers}. 
%We observe the trend that $\lambda = 20d_0$ works best in stabilizing the PDEs for all cases. 

Recall that the control objective \eqref{eq:L2objective} in finite time is to drive $\|z(T,\cdot)\|_{L^2}$  as close to zero as possible. To illustrate the performance of the controllers we make two plots for both PDEs. 
% The first plot is a heat map depicting the mean value (over the 100 simulations) of $\frac{\|z(T,\cdot)\|_{L^2}}{\|z(0,\cdot)\|_{L^2}}$ -- the ratio of L2 norm at the final time to L2 norm at initial time. The ratio is plotted to show the reduction in energy of the PDE relative to its initial energy. 
The first plot is a heat map depicting the mean value (over the 100 simulations) of $\frac{\|z(T,\cdot)\|_{L^2}}{\|z(0,\cdot)\|_{L^2}}$ -- the ratio is plotted to highlight the order of magnitude by which $\|z(T,\cdot)\|_{L^2}$  has reduced relative to its initial value $\|z(0,\cdot)\|_{L^2}$. 
Moreover, in a second plot, we plot the mean and variance (over the 100 simulations) of  $\|z(t,\cdot)\|_{L^2}$ as a function of $t$ for $d_0 = 0.1$ and $\lambda = 0.2$. The trajectory is generated using three different control policies -- ``uncontrolled" trajectories are when the control is zero, ``optimal controlled" trajectories have $\lambda = 0$ (that is, only the optimal stabilizing control is applied and robust control is zero) and ``robust controlled" trajectories are with $\lambda = 0.2$. The plots for heat equation, Burgers with DMDc and Burgers with full nonlinear model are found in Figures \ref{fig:heat}, \ref{fig:dmdc}, and \ref{fig:full}, respectively. Other simulations details can be obtained in Appendix \ref{app:burgers}. 

For both PDEs, we observe that in the presence of disturbances, the robust control works better than the optimal control (control with $\lambda = 0$) in stabilizing the PDE. For the Burgers' equation in particular, the trajectories with robust control exhibit an order of magnitude lower value of $\|z(T,\cdot)\|_{L^2}$ compared to those with only the optimal control. Additionally, for the Burgers' equation, we also observe that the optimal control obtained using the full nonlinear model is significantly more effective than using the reduced-order DMDc model -- the controlled state settles much faster and closer to zero with the former than the latter. However, while the transient performance of the robust control is much better using the full nonlinear model, the settling performance is similar using both DMDc and the full nonlinear model. Thus the robust control term also compensates for model mismatch between DMDc and the full nonlinear model. The conclusions presented are true for all values of $\lambda$ chosen, both types of disturbances, and both values of viscosity. 
% \todo{what is exactly plotted? also include discussion on the results, see comments below}




% \begin{figure*}
% \centering
%     \begin{subfigure}[b]{0.45\textwidth}
%         \centering
%         \includegraphics[scale = 0.35]{Fig/heatmap-dmdc.pdf}
%         \caption{dmdc}
%         \label{fig:heatmap-dmdc}
%     \end{subfigure}
% %
%     \begin{subfigure}[b]{0.45\textwidth}
%         \centering
%         \includegraphics[scale = 0.35]{Fig/heatmap-full.pdf}
%         \caption{full}
%         \label{fig:heatmap-full}
%     \end{subfigure}
% \caption{hwfk}
% \label{fig:heatmap-burgers}
% \end{figure*}

% \begin{figure*}
% \centering
%     \begin{subfigure}[b]{0.45\textwidth}
%         \centering
%         \includegraphics[scale = 0.35]{Fig/time-dmdc.pdf}
%         \caption{dmdc}
%         \label{fig:heatmap-dmdc}
%     \end{subfigure}
% %
%     \begin{subfigure}[b]{0.45\textwidth}
%         \centering
%         \includegraphics[scale = 0.35]{Fig/time-full.pdf}
%         \caption{full}
%         \label{fig:heatmap-full}
%     \end{subfigure}
% \caption{hwfk}
% \label{fig:time-burgers}
% \end{figure*}

\begin{figure*}
\centering
    \begin{subfigure}[b]{0.45\textwidth}
        \centering
        
        \caption{Mean of $\frac{\|z(T,\cdot)\|_{L^2}}{\|z(0,\cdot)\|_{L^2}}$ over 100 different simulations}
        \label{fig:heatmap-dmdc}
    \end{subfigure}
%
    \begin{subfigure}[b]{0.45\textwidth}
        \centering
        
        \caption{Mean and variance of $\|z(t,\cdot)\|_{L^2}$ for $d_0=0.1,\lambda=0.2$.}
        \label{fig:time-dmdc}
    \end{subfigure}
\caption{Results for control of Burgers equation using reduced order DMDc model.}
\label{fig:dmdc}
\end{figure*}

\begin{figure*}
\centering
    \begin{subfigure}[b]{0.45\textwidth}
        \centering
        
        \caption{Mean of $\frac{\|z(T,\cdot)\|_{L^2}}{\|z(0,\cdot)\|_{L^2}}$ over 100 different simulations}
        \label{fig:heatmap-full}
    \end{subfigure}
%
    \begin{subfigure}[b]{0.45\textwidth}
        \centering
        
        \caption{Mean and variance of $\|z(t,\cdot)\|_{L^2}$ for $d_0=0.1,\lambda=0.2$.}
        \label{fig:time-full}
    \end{subfigure}
\caption{Results for control of Burgers equation using full nonlinear model.}
\label{fig:full}
\end{figure*}

% \begin{figure}
%     \centering
%     \includegraphics[scale = 0.4]{Fig/heatmap-full.pdf}
%     \caption{Caption}
%     \label{fig:heatmap-full}
% \end{figure}

\section{Future work}
Some avenues for future work are: (i) to extend this approach to different PDEs such as the Allen-Cahn equation or the Korteweg deVries equation, (ii) to design a controller using information from sensors, thus turning a full state feedback problem into a partially observed problem (iii) to extend to the case when $d$ is Gaussian white noise. 
%Another major avenue for future work is to improve the robust control design procedure by choosing more appropriate functions for $\lambda$ to account for model mismatch.

\section{Dual EnKF algorithm}
\label{app:enkf}
In this algorithm, we simulate over the time horizon $[0,T]$ an ensemble of $N$ particles $\{ Y_t^i : 1 \le i \le N, \quad 0 \le t \le T \}$ where the evolution equation for each particle is given by an Ito stochastic differential equation (SDE). 

\subsection{Linear system}
\label{app:linenkf}
The SDE \cite[Section 2]{joshi-2022} is given by
%\begin{subequations}\label{eq:EnKF-Yi}
\begin{align*}
\ud {Y}^i_t &= {A {Y}^i_t \ud t + B\ud \backward{\eta}^i_t } + 
L^{(N)}_t \left(\frac{C{Y}^i_t+ C\hat{n}^{(N)}_t}{2} \right), \\ 
{Y}^i_T  &\stackrel{\text{i.i.d}}{\sim} \mathcal
N(0,S_T),\quad 1\leq i\leq N, \\
\end{align*}
%\end{subequations}
${\eta}^i$ are an i.i.d Brownian motions with covariance $R\inv$ and $\hat{n}^{(N)}_t := \left( \frac{1}{N} \sum_{j=1}^{N} Y_t^j \right) $, and $L_t^{(N)} := S_t^{(N)}C\tp$ with
\begin{align*}
S^{(N)}_t := \frac{1}{N} \sum_{j=1}^{N} ({Y}_t^j - \hat{n}_t^{(N)}) ({Y}_t^j - \hat{n}_t^{(N)})\tp. 
\end{align*}
Finally, $\bar{P}^{(N)} := (S^{(N)}_0)\inv$.

\subsection{Nonlinear system}
\label{app:nlenkf}
The SDE \cite[Section 3]{joshi-2022} for the particle system is
%\begin{subequations}
\begin{align*}
\ud {Y}^i_t &= {a ({ Y}^i_t)\ud t +  b ({ Y}^i_t) \ud \backward{\eta}^i_t}  
 \\
 &+ \left(\sum_{j=1}^{N}(Y_t^j - n^{(N)}_t )(c({ Y}^j_t) - \hat{c}_t^{(N)}  )^\top\right) \frac{(c(z) + \hat{c}^{(N)})}{N-1} \\
{Y}^i_T  &\stackrel{\text{i.i.d}}{\sim} \mathcal
N(0,S_T),\quad 1\leq i\leq N, 
\end{align*}
%\end{subequations}
and $\hat{c}^{(N)}_{t} := N^{-1} \sum_{i=1}^N c(Y^i_{t})$. Finally, $\grad\phi^{(N)}(x) := (S^{(N)}_0)\inv x$, where $S^{(N)}_t$ is the empirical covariance of $\{Y_t^i\}$, defined same as in the linear case. 
%and the gain is a constant matrix:
% \[
% {\sf K}_t^{(N)} = \frac{1}{N-1}\sum_{i=1}^N  ({ Y}^i_t-n^{(N)}_t )(c({ Y}^i_t) - \hat{c}_t^{(N)}  )^\top 
% \]
%One may interpret the above as the dual counterpart of the FPF
%algorithm with a constant gain approximation~\cite[Example 2]{yang-2016}. 

\subsection{Algorithm for Hamiltonian minimization }
\label{app:ubar}
Algorithm \ref{alg:ubar} minimizes Hamiltonian and calculates $\bar{u}$ for Section \ref{sec:linear}. To minimize Hamiltonian in Section \ref{sec:nonlinear}, in Algorithm \ref{alg:ubar} make the following two changes: replace $\bar{P}^{(N)}$  by $\grad\phi^{(N)}$ and use the Hamiltonian defined in \eqref{eq:ham-nl}.
\begin{algorithm}
\begin{algorithmic}[1]
\Input System state $x$, $\bar{P}^{(N)}$, $Q,R$, Hamiltonian definition from \eqref{eq:ham-lin}, $\{e_i\}_{i=1}^{m}$ the standard basis of $\R^m$
\If{$B$ is known}
\State \Return $\bar{u}^{(N)} := -R\inv B\tp \bar{P}^{(N)}x$
\ElsIf{$B$ is unknown}
\For{i=1,2,\ldots,m}
    \State $(\bar{u}^{(N)})_i := H(x,R\inv e_i) - H(x,0) - \half (R\inv)_{ii}$
    \EndFor
\EndIf
\end{algorithmic}
\caption{Algorithm for Hamiltonian minimization}
\label{alg:ubar}
\end{algorithm}

\section{Simulation details}
\subsection{Heat equation}
\label{app:heat}
Simulation parameters are as follows. The simulation time $T = 0.1$ %and 10
with simulation time step = 0.001. 
The number of states of the discretized PDE \eqref{eq:discrete-pde} is $p=100$. 
The number of control basis functions is $m=8$, and they are $\chi_j$ is the indicator functions of $[\frac{i}{10},\frac{(i+1)}{10}]$. The regularization parameter for robust control is $r = 0.002$.
The matrices  $Q = \id$, $G=\id$, $R=\id$ for both, the full nonlinear control and control using DMDc model.
The number of dual EnKF particles is chosen as $N=10000$. The controlgym library \cite{zhang2023controlgym} is used for PDE simulation. 

\subsection{Burgers equation}
\label{app:burgers}
Simulation parameters are as follows. The simulation time $T = 3$ %and 10
with simulation time step = 0.001. 
The number of states of the discretized PDE \eqref{eq:discrete-pde} is $p=128$. 
The number of control basis functions is $m=10$, and they are $\chi_j$ is the indicator functions of $[\frac{i}{10},\frac{(i+1)}{10}]$. The regularization parameter for robust control is $r = 0.002$.
The number of states in the DMDc model \eqref{eq:DMDc} is $n=10$. The matrices  $Q = \id$, $G=\id$, $R=0.1\id$ for both, the full nonlinear control and control using DMDc model.
The number of dual EnKF particles is chosen as $N=1000$. 
%\todo{Mention number of particles, dimension of state and control, value of $Q,R,G$, sim time step, distribution of initial states, r, sim time T for EnKF and PDE}
The controlgym library \cite{zhang2023controlgym} is used for PDE simulation. 



\begin{thebibliography}{99}
\bibitem[1]{bach-2016} {\sc F.~Bach and V.~Perchet}, {\em Highly-smooth zero-th order online optimization}, in 29th Annual Conference on Learning Theory, V.~Feldman, A.~Rakhlin, and O.~Shamir, eds., vol.~49 of Proceedings of Machine Learning Research, Columbia University, New York, New York, USA, 23--26 Jun 2016, PMLR, pp.~257--283.
\bibitem[2]{barbagallo2009closed} {\sc A.~Barbagallo, D.~Sipp, and P.~J. Schmid}, {\em Closed-loop control of an open cavity flow using reduced-order models}, Journal of Fluid Mechanics, 641 (2009), pp.~1--50.
\bibitem[3]{bhattacharya2021model} {\sc K.~Bhattacharya, B.~Hosseini, N.~B. Kovachki, and A.~M. Stuart}, {\em Model reduction and neural networks for parametric {PDE}s}, The SMAI Journal of Computational Mathematics, 7 (2021), pp.~121--157.
\bibitem[4]{bishop-2023} {\sc A.~N. Bishop and P.~Del~Moral}, {\em On the mathematical theory of ensemble (linear-gaussian) kalman--bucy filtering}, Mathematics of Control, Signals, and Systems, 35 (2023), pp.~835--903.
\bibitem[5]{blanchard2021bayesian} {\sc A.~B. Blanchard, G.~Y. Cornejo~Maceda, D.~Fan, Y.~Li, Y.~Zhou, B.~R. Noack, and T.~P. Sapsis}, {\em Bayesian optimization for active flow control}, Acta Mechanica Sinica, (2021), pp.~1--13.
\bibitem[6]{duriez2017machine} {\sc T.~Duriez, S.~L. Brunton, and B.~R. Noack}, {\em Machine Learning Control-Taming Nonlinear Dynamics and Turbulence}, vol.~116, Springer, 2017.
\bibitem[7]{evensen2006} {\sc G.~Evensen}, {\em Data Assimilation. {T}he Ensemble {K}alman Filter}, Springer-Verlag, New York, 2006.
\bibitem[8]{fan2020reinforcement} {\sc D.~Fan, L.~Yang, Z.~Wang, M.~S. Triantafyllou, and G.~E. Karniadakis}, {\em Reinforcement learning for bluff body active flow control in experiments and simulations}, Proceedings of the National Academy of Sciences, 117 (2020), pp.~26091--26098.
\bibitem[9]{fazel-2018-mlr} {\sc M.~Fazel, R.~Ge, S.~Kakade, and M.~Mesbahi}, {\em Global convergence of policy gradient methods for the linear quadratic regulator}, in Proceedings of the 35th International Conference on Machine Learning, J.~Dy and A.~Krause, eds., vol.~80 of Proceedings of Machine Learning Research, PMLR, 10--15 Jul 2018, pp.~1467--1476.
\bibitem[10]{fleming-1982} {\sc W.~H. Fleming and S.~K. Mitter}, {\em Optimal {Control} and {Nonlinear} {Filtering} for {Nondegenerate} {Diffusion} {Processes}}, Stochastics, 8 (1982), pp.~63--77.
\bibitem[11]{fresca2021comprehensive} {\sc S.~Fresca, L.~Dede’, and A.~Manzoni}, {\em A comprehensive deep learning-based approach to reduced order modeling of nonlinear time-dependent parametrized {PDE}s}, Journal of Scientific Computing, 87 (2021), pp.~1--36.
\bibitem[12]{Garnieretal21} {\sc P.~Garnier, J.~Viquerat, J.~Rabault, A.~Larcher, A.~Kuhnle, and E.~Hachem}, {\em A review on deep reinforcement learning for fluid mechanics}, Computers \& Fluids, 225 (2021), p.~104973.
\bibitem[13]{hoffman-2017} {\sc C.~Hoffmann and P.~Rostalski}, {\em Linear optimal control on factor graphs — a message passing perspective —}, IFAC-PapersOnLine, 50 (2017), pp.~6314--6319. \newblock 20th IFAC World Congress.
\bibitem[14]{hovland2008explicit} {\sc S.~Hovland, J.~T. Gravdahl, and K.~E. Willcox}, {\em Explicit model predictive control for large-scale systems via model reduction}, Journal of guidance, control, and dynamics, 31 (2008), pp.~918--926.
\bibitem[15]{joshi-2022} {\sc A.~A. Joshi, A.~Taghvaei, P.~G. Mehta, and S.~P. Meyn}, {\em Controlled interacting particle algorithms for simulation-based reinforcement learning}, Systems \& Control Letters, 170 (2022), p.~105392.
\bibitem[16]{kaiser2021data} {\sc E.~Kaiser, J.~N. Kutz, and S.~L. Brunton}, {\em Data-driven discovery of {K}oopman eigenfunctions for control}, Machine Learning: Science and Technology, 2 (2021), p.~035023.
\bibitem[17]{kappen-2005} {\sc H.~J. Kappen}, {\em Linear theory for control of nonlinear stochastic systems}, Phys. Rev. Lett., 95 (2005), p.~200201.
\bibitem[18]{khalil} {\sc H.~K. Khalil}, {\em Nonlinear systems}, Macmillan Pub. Co., New York, 1992.
\bibitem[19]{k-sivan} {\sc H.~Kwakernaak and R.~Sivan}, {\em Linear optimal control systems}, Wiley Interscience, New York, 1972.
\bibitem[20]{leibfritz2007numerical} {\sc F.~Leibfritz and S.~Volkwein}, {\em Numerical feedback controller design for {PDE} systems using model reduction: {T}echniques and case studies}, in Real-Time PDE-Constrained Optimization, SIAM, 2007, pp.~53--72.
\bibitem[21]{levine-2018} {\sc S.~Levine}, {\em Reinforcement learning and control as probabilistic inference: Tutorial and review}, 2018.
\bibitem[22]{liberzon} {\sc D.~Liberzon}, {\em Calculus of variations and optimal control theory}, Princeton University Press, Princeton, NJ, 2012.
\bibitem[23]{mihailo-2021-tac} {\sc H.~Mohammadi, A.~Zare, M.~Soltanolkotabi, and M.~R. Jovanović}, {\em Convergence and sample complexity of gradient methods for the model-free linear–quadratic regulator problem}, IEEE Transactions on Automatic Control, 67 (2022), pp.~2435--2450.
\bibitem[24]{proctor2016dynamic} {\sc J.~L. Proctor, S.~L. Brunton, and J.~N. Kutz}, {\em Dynamic mode decomposition with control}, SIAM Journal on Applied Dynamical Systems, 15 (2016), pp.~142--161.
\bibitem[25]{sipp2016linear} {\sc D.~Sipp and P.~J. Schmid}, {\em Linear closed-loop control of fluid instabilities and noise-induced perturbations: A review of approaches and tools}, Applied Mechanics Reviews, 68 (2016), p.~020801.
\bibitem[26]{amir-2019} {\sc A.~Taghvaei, J.~de~Wiljes, P.~G. Mehta, and S.~Reich}, {\em Kalman filter and its modern extensions for the continuous-time nonlinear filtering problem}, Journal of Dynamic Systems, Measurement, and Control, 140 (2017), p.~030904.
\bibitem[27]{todorov-2007} {\sc E.~Todorov}, {\em Linearly-solvable markov decision problems}, in Advances in Neural Information Processing Systems, B.~Sch\"{o}lkopf, J.~Platt, and T.~Hoffman, eds., vol.~19, MIT Press, 2007.
\bibitem[28]{tsolovikos2020estimation} {\sc A.~Tsolovikos, E.~Bakolas, S.~Suryanarayanan, and D.~Goldstein}, {\em Estimation and control of fluid flows using sparsity-promoting dynamic mode decomposition}, IEEE Control Systems Letters, 5 (2020), pp.~1145--1150.
\bibitem[29]{rawlik2013stochastic} {\sc S.~Vijayakumar, K.~Rawlik, and M.~Toussaint}, {\em On stochastic optimal control and reinforcement learning by approximate inference}, in Robotics: Science and Systems VIII, N.~Roy, P.~Newman, and S.~Srinivasa, eds., MIT Press, 2013, pp.~353--360.
\bibitem[30]{yang-2016} {\sc T.~Yang, R.~S. Laugesen, P.~G. Mehta, and S.~P. Meyn}, {\em Multivariable feedback particle filter}, Automatica, 71 (2016), pp.~10--23.
\bibitem[31]{zhang2023controlgym} {\sc X.~Zhang, W.~Mao, S.~Mowlavi, M.~Benosman, and T.~Ba{\c{s}}ar}, {\em Controlgym: Large-scale safety-critical control environments for benchmarking reinforcement learning algorithms}, arXiv preprint arXiv:2311.18736, (2024).
\bibitem[32]{zhang-2024} {\sc X.~Zhang, S.~Mowlavi, M.~Benosman, and T.~Başar}, {\em Policy optimization for pde control with a warm start}, arXiv preprint arXiv 2403.01005, (2024).
\end{thebibliography}
\end{document}
